\clearpage
\SourceLeaf{106}\pdfbookmark[1]{Section III.}{sectio-tertia}
\begin{center}
{\large SECTION III.}\par\medskip
\emph{Relations among several positions in the orbit.}\par\bigskip
\hypertarget{native-106-article078}{}\pdfbookmark[2]{78.}{article-078}78.
\end{center}
Comparing two or more positions of a celestial body, both in its orbit and in space, yields so many elegant propositions that they could easily fill an entire volume. An exhaustive treatment of this rich subject lies outside our purpose. Our chief aim is to draw from it ample means for solving the important problem of determining unknown orbits from observations. We therefore leave aside what lies too far from our purpose and develop with all the more care everything that can contribute to it in any way. Before beginning the investigation itself, we give some trigonometric propositions that are frequently useful and to which we shall have to return repeatedly.
\par
I. Let $A$, $B$, and $C$ be arbitrary angles. Then
\[
\begin{aligned}
\sin A\sin(C-B)+\sin B\sin(A-C)+\sin C\sin(B-A)&=0\\
\cos A\sin(C-B)+\cos B\sin(A-C)+\cos C\sin(B-A)&=0.
\end{aligned}
\]
\par
II. Suppose two quantities $p$ and $P$ are to be determined from equations such as
\[
\begin{aligned}
p\sin(A-P)&=a\\
p\sin(B-P)&=b
\end{aligned}
\]
In general this is done by means of the formulas
\SourceLeaf{107}
\[
\begin{aligned}
p\sin(B-A)\sin(H-P)&=b\sin(H-A)-a\sin(H-B)\\
p\sin(B-A)\cos(H-P)&=b\cos(H-A)-a\cos(H-B),
\end{aligned}
\]
where $H$ is an arbitrary angle. From these formulas we obtain (Article 14, II) the angle $H-P$ and $p\sin(B-A)$, and from these, $P$ and $p$. It is usual to require that $p$ be positive. This resolves the ambiguity in determining $H-P$ from its tangent. In the absence of that condition, either choice is permissible. To make the calculation as convenient as possible, we take the arbitrary angle $H$ to be $=A$, $=B$, or $=\tfrac{1}{2}(A+B)$. In the first case, the equations for determining $P$ and $p$ are
\[
\begin{aligned}
p\sin(A-P)&=a\\
p\cos(A-P)&=\frac{b-a\cos(B-A)}{\sin(B-A)}\mathpunct{\cdot}
\end{aligned}
\]
The equations in the second case are entirely analogous. In the third case they are
\[
\begin{aligned}
p\sin(\tfrac{1}{2}A+\tfrac{1}{2}B-P)&=\frac{b+a}{2\cos\tfrac{1}{2}(B-A)}\\
p\cos(\tfrac{1}{2}A+\tfrac{1}{2}B-P)&=\frac{b-a}{2\sin\tfrac{1}{2}(B-A)}\mathpunct{\cdot}
\end{aligned}
\]
\par
If we therefore introduce the auxiliary angle $\zeta$, whose tangent is $=\frac{a}{b}$, we obtain $P$ by the formula
\[
\operatorname{tang}(\tfrac{1}{2}A+\tfrac{1}{2}B-P)=\operatorname{tang}(45^0+\zeta)\operatorname{tang}\tfrac{1}{2}(B-A)
\]
and then $p$ by one of the preceding formulas, using
\[
\begin{aligned}
\tfrac{1}{2}(b+a)&=\sin(45^0+\zeta)\BareRoot{\frac{ab}{\sin2\zeta}}\\
&=\frac{a\sin(45^0+\zeta)}{\sin\zeta\BareRoot{2}}=\frac{b\sin(45^0+\zeta)}{\cos\zeta\BareRoot{2}}\\[.4em]
\tfrac{1}{2}(b-a)&=\cos(45^0+\zeta)\BareRoot{\frac{ab}{\sin2\zeta}}\\
&=\frac{a\cos(45^0+\zeta)}{\sin\zeta\BareRoot{2}}=\frac{b\cos(45^0+\zeta)}{\cos\zeta\BareRoot{2}}\mathpunct{\cdot}
\end{aligned}
\]
\par
III. Suppose $p$ and $P$ are to be determined from the equations
\[
\begin{aligned}
p\cos(A-P)&=a\\
p\cos(B-P)&=b,
\end{aligned}
\]
Everything set out in II would apply at once if $A$ and $B$ were everywhere replaced there by $90^0+A$ and $90^0+B$. For greater convenience in using them, we give the resulting formulas explicitly. The general formulas are
\SourceLeaf{108}
\[
\begin{aligned}
p\sin(B-A)\sin(H-P)&=-b\cos(H-A)+a\cos(H-B)\\
p\sin(B-A)\cos(H-P)&=b\sin(H-A)-a\sin(H-B).
\end{aligned}
\]
For $H=A$ these thus become
\[
\begin{aligned}
p\sin(A-P)&=\frac{a\cos(B-A)-b}{\sin(B-A)}\\
p\cos(A-P)&=a.
\end{aligned}
\]
For $H=B$ they take a similar form. For $H=\tfrac{1}{2}(A+B)$ they become
\[
\begin{aligned}
p\sin(\tfrac{1}{2}A+\tfrac{1}{2}B-P)&=\frac{a-b}{2\sin\tfrac{1}{2}(B-A)}\\
p\cos(\tfrac{1}{2}A+\tfrac{1}{2}B-P)&=\frac{a+b}{2\cos\tfrac{1}{2}(B-A)}\mathpunct{\raisebox{.45ex}{,}}
\end{aligned}
\]
so that on introducing the auxiliary angle $\zeta$, whose tangent is $=\frac{a}{b}$, we have
\[
\operatorname{cotang}(\tfrac{1}{2}A+\tfrac{1}{2}B-P)=\operatorname{cotang}(\zeta-45^0)\operatorname{tang}\tfrac{1}{2}(B-A).
\]
\par
If, moreover, we wish to determine $p$ directly from $a$ and $b$ without first calculating the angle $P$, we have the formula
\[
p\sin(B-A)=\BareRoot{(aa+bb-2ab\cos(B-A))}
\]
both in the present problem and in II.
